\subsection{Дизайн-освіта в Україні: розвиток компетентностей та європейська інтеграція}\label{subsec:1.3}

Попередній підрозділ засвідчив, що трансформація дизайн-освіти під впливом генеративного ШІ потребує переосмислення компетентнісних вимог до фахівців-дизайнерів. Однак ця трансформація відбувається не в ізольованому просторі, а в конкретному інституційному, нормативному та геополітичному контексті. Для України, яка з 2005~р. є учасницею Болонського процесу, ключовим чинником такого контексту є процес європейської інтеграції вищої освіти й необхідність узгодження національних стандартів з європейськими рамками кваліфікацій. У цьому підрозділі здійснюється систематичний аналіз компетентнісних структур українського стандарту вищої освіти B2 <<Дизайн>> у зіставленні з Європейською рамкою кваліфікацій (ЄРК), досліджуються педагогічні підходи до формування художньо-проєктних компетентностей та визначаються виклики і можливості європейської інтеграції для дизайн-освіти в Україні. Особливу увагу приділено зв'язку між компетентностями стандарту B2 та компонентами інформаційно-комунікативної компетентності (ІКК) як інтегративної характеристики підготовки майбутніх дизайнерів.

Основу аналізу складає теоретичне дослідження \cite{Musiienko2025DesignEducation}, у якому автором було здійснено порівняльний аналіз українського стандарту дизайн-освіти, європейських кваліфікаційних рамок та міжнародних моделей дизайн-освіти. Нижче результати цього дослідження адаптовано, розширено та переосмислено в контексті формування ІКК.

\textbf{Стандарт B2 <<Дизайн>>: компетентнісна структура.} Стандарт вищої освіти України за спеціальністю 022 <<Дизайн>> (бакалаврський рівень), затверджений Міністерством освіти і науки у 2018~р.~\cite{MON2018}, визначає компетентнісні вимоги до підготовки фахівців-дизайнерів, що мають забезпечити їхню здатність розв'язувати складні дизайнерські задачі в умовах неповної інформації. Стандарт організовує компетентності у дві категорії: загальні компетентності (ЗК) та фахові компетентності (ФК).

Загальні компетентності (ЗК1--ЗК14) охоплюють критичне мислення та аналіз (ЗК1), здатність до пошуку, оброблення та аналізу інформації з різних джерел (ЗК7), навички використання інформаційних та комунікаційних технологій (ЗК8), здатність спілкуватися іноземною мовою (ЗК5), а також уміння працювати в команді (ЗК6) та спілкуватися з представниками інших професійних груп різного рівня (ЗК9). Фахові компетентності (ФК1--ФК14) зосереджені на здатності здійснювати художньо-проєктну діяльність (ФК1), проводити естетичний аналіз об'єктів дизайну (ФК3), застосовувати сучасне програмне забезпечення для створення дизайн-об'єктів (ФК9), а також здатності до художньо-образного мислення (ФК4).

Зазначена структура безпосередньо перегукується зі структурою ІКК, визначеною в підрозділі~\ref{subsec:1.1}. Зокрема, загальні компетентності ЗК7 та ЗК8 формують базис інформаційно-аналітичного компонента ІКК (пошук, оцінювання, використання професійної інформації); ЗК5, ЗК6 та ЗК9 є фундаментом комунікативного компонента (цифрова професійна комунікація, візуальна презентація, колаборація); ФК9 корелює з технологічним компонентом (використання цифрових інструментів дизайну); а ЗК1 разом з ФК3 та ФК4 утворюють основу рефлексивного компонента (критичне оцінювання, естетичний аналіз, самооцінка). Тим самим стандарт B2 імпліцитно визначає структурні елементи ІКК, хоча й не артикулює їх як цілісну інтегративну характеристику фахівця.

\textbf{Зіставлення зі стандартом ЄРК: структурні відповідності.} Європейська рамка кваліфікацій встановлює тричастинну категоризацію результатів навчання --- знання, уміння та компетентності (відповідальність і автономія), --- яка є нормативною для розроблення навчальних програм у межах Європейського простору вищої освіти~\cite{EuropeanCommission2017}. Стандарт B2 демонструє чітку структурну відповідність архітектурі ЄРК: поділ компетентностей на загальні (ЗК) та фахові (ФК) відображає розмежування між трансверсальними та дисциплінарно-специфічними здатностями. Обидві рамки позиціонують бакалаврську кваліфікацію на рівні~6, що передбачає продемонстровану здатність застосовувати знання у складних, непередбачуваних контекстах із відповідальністю за прийняття рішень~\cite{Musiienko2025DesignEducation}.

Порівняльний аналіз виявляє сильну відповідність за вимірами знань та умінь: обидві рамки акцентують розвинені спеціалізовані здатності, що застосовуються до розв'язання складних проблем~\cite{Bohlinger2012279}. Наголос стандарту B2 на художньо-проєктній діяльності як інтегрувальному домені компетентностей добре кореспондує з очікуваннями ЄРК щодо дисциплінарної майстерності, продемонстрованої через практичне застосування. У вимірі відповідальності та автономії спостерігається помірна відповідність: ЄРК більш артикульовано визначає управління складною технічною або професійною діяльністю, тоді як B2 формулює це опосередковано через здатність працювати самостійно та в команді.

Водночас існують суттєві термінологічні варіації, що можуть ускладнювати міжрамкову інтерпретацію. Стандарт B2 оперує концептом \emph{художньо-проєктної діяльності} (ХПД) як центральним організаційним принципом фахових компетентностей. Цей концепт, укорінений у радянській традиції дизайн-освіти, інтегрує художню та технічну складові у спосіб, що не має прямих еквівалентів у західноєвропейських рамках~\cite{Fursa2014}. Тричастинна структура ЄРК (знання --- уміння --- компетентності) відрізняється від двочастинної організації стандарту B2 (загальні та фахові компетентності). Хоча ці структури можна зіставити, така структурна відмінність відображає глибші концептуальні варіації: ЄРК трактує компетентність як окрему категорію, що охоплює відповідальність та автономію, тоді як B2 розподіляє ці елементи між двома категоріями компетентностей~\cite{Gehmlich2009736}.

Крім того, фахові компетентності стандарту B2 акцентують естетичні та художні виміри значно більш артикульовано, ніж типові рамки, узгоджені з ЄРК, для технічних дисциплін. Такі компетентності як <<здатність проводити естетичний аналіз>> (ФК3) та <<здатність до художньо-образного мислення>> (ФК4) висувають на передній план креативні здатності, які дескриптори ЄРК адресують лише імпліцитно через загальні посилання на інновацію та розв'язання проблем~\cite{Musiienko2025DesignEducation}.

Результати зіставлення узагальнено в табл.~\ref{tab:b2-eqf-alignment}, яка демонструє як структурні відповідності, так і наявні розбіжності.

\begin{table}[htbp]
\caption{Зіставлення компетентностей стандарту B2 <<Дизайн>> та дескрипторів ЄРК рівня~6}
\label{tab:b2-eqf-alignment}
\centering\small
\begin{tabularx}{\textwidth}{@{}p{2.5cm}XXc@{}}
\toprule
\textbf{Вимір} & \textbf{Дескриптор ЄРК рівня~6} & \textbf{Специфікація стандарту B2} & \textbf{Ступінь} \\
\midrule
Знання & Поглиблені знання у галузі роботи чи навчання, включно з критичним розумінням теорій та принципів & Знання з історії, теорії та методології дизайну; розуміння принципів художньо-проєктної діяльності & Сильний \\
\addlinespace
Уміння & Поглиблені вміння, що демонструють майстерність та інноваційність у розв'язанні складних і непередбачуваних проблем & Уміння у сфері ХПД; здатність розв'язувати складні дизайнерські задачі з неповною інформацією & Сильний \\
\addlinespace
Відповідальність та автономія & Управління складною технічною або професійною діяльністю; відповідальність за прийняття рішень & Здатність працювати самостійно та в команді; відповідальність за професійний розвиток & Помірний \\
\addlinespace
Критичне мислення & Імпліцитно --- через розв'язання проблем та прийняття рішень & ЗК1: критичне мислення та аналіз & Сильний \\
\addlinespace
Комунікація & Імпліцитно --- в професійних контекстах & ЗК5, ЗК6, ЗК9: професійна комунікація, командна робота & Сильний \\
\addlinespace
Інформаційні технології & Імпліцитно --- через загальні вимоги до рівня~6 & ЗК7, ЗК8: пошук інформації, ІКТ-навички; ФК9: програмне забезпечення дизайну & Сильний \\
\bottomrule
\end{tabularx}
\end{table}

\textbf{Зв'язок компетентностей B2 з компонентами ІКК.} Наведений вище компетентнісний аналіз дозволяє встановити систематичні кореляції між компетентностями стандарту B2 та чотирикомпонентною структурою ІКК. Таке зіставлення має принципове значення, оскільки демонструє, що ІКК не є зовнішнім доповненням до стандарту, а являє собою інтегративну характеристику, яка синтезує декілька нормативно визначених компетентностей у єдину функціональну цілісність (табл.~\ref{tab:b2-ikk-mapping}).

\begin{table}[htbp]
\caption{Зіставлення компетентностей стандарту B2 <<Дизайн>> з компонентами ІКК}
\label{tab:b2-ikk-mapping}
\centering\small
\begin{tabularx}{\textwidth}{@{}p{2.8cm}p{3.2cm}Xp{2.5cm}@{}}
\toprule
\textbf{Компонент ІКК} & \textbf{Компетентності B2} & \textbf{Зміст зв'язку} & \textbf{Відповідність ЄРК} \\
\midrule
Інформаційно-аналітичний & ЗК1 (критичне мислення), ЗК7 (пошук інформації) & Пошук, критичне оцінювання та використання професійної інформації, зокрема генерованої ШІ & Знання; критичне мислення \\
\addlinespace
Комунікативний & ЗК5 (іноземна мова), ЗК6 (командна робота), ЗК9 (міжгрупова комунікація) & Цифрова професійна комунікація, візуальна презентація, міждисциплінарна колаборація & Комунікація; відповідальність \\
\addlinespace
Технологічний & ЗК8 (ІКТ-навички), ФК9 (програмне забезпечення дизайну) & Цифрові інструменти дизайну (традиційні та ШІ), промпт-інженерія, адаптивний вибір засобів & Уміння (технічні) \\
\addlinespace
Рефлексивний & ЗК1 (критичне мислення), ФК3 (естетичний аналіз), ФК4 (художньо-образне мислення) & Самооцінка, критичне оцінювання ШІ в дизайні, етика, безперервне навчання & Відповідальність та автономія \\
\bottomrule
\end{tabularx}
\end{table}

Як видно з табл.~\ref{tab:b2-ikk-mapping}, кожен компонент ІКК забезпечується щонайменше двома компетентностями стандарту B2, що підтверджує інтегративний характер ІКК. При цьому варто зазначити, що стандарт B2 не містить прямих посилань на ШІ-грамотність, промпт-інженерію або етику штучного інтелекту, що створює лакуну у нормативному забезпеченні формування ІКК, релевантної сучасним вимогам дизайн-практики. Ця лакуна кореспондує з виявленим у попередньому підрозділі переходом від парадигми <<дизайнер як виробник>> до парадигми <<дизайнер як куратор>>, яка суттєво змінює зміст інформаційно-комунікативної діяльності фахівця.

\textbf{Виявлені компетентнісні лакуни.} Порівняльний аналіз стандарту B2 з сучасними міжнародними моделями дизайн-освіти дозволяє ідентифікувати три ключові сфери, в яких український стандарт виявляє менш розвинену відповідність~\cite{Musiienko2025DesignEducation}.

\emph{По-перше, компетентності критичної рефлексії.} Стандарт B2 включає критичне мислення як загальну компетентність (ЗК1), однак артикульована увага до рефлексивної практики --- систематичного осмислення власного дизайнерського процесу та рішень --- виявляється недорозвиненою порівняно з міжнародними моделями, які висувають на передній план метакогнітивні здатності~\cite{Schon1983, Hameed2023310}. Це є суттєвим для формування рефлексивного компонента ІКК, оскільки саме рефлексивна практика забезпечує здатність критично оцінювати результати взаємодії з ШІ-системами, усвідомлювати обмеження автоматизованих рішень та підтримувати авторську позицію в дизайнерському процесі.

\emph{По-друге, міждисциплінарна колаборація.} Сучасний дизайн дедалі більше функціонує на перетині з інженерією, бізнесом, соціальними науками та іншими галузями~\cite{Davis202397}. Стандарт B2 адресує командну роботу на загальному рівні (ЗК6), проте не містить специфікації компетентностей міждисциплінарної колаборації, які набувають центрального значення у професійній практиці дизайнера. Для комунікативного компонента ІКК це означає необхідність розширення за межі суто дизайнерської комунікації до здатності вести продуктивний діалог з фахівцями інших галузей, зокрема з розробниками ШІ-систем, фахівцями з обробки даних та UX-дослідниками.

\emph{По-третє, сталий розвиток та соціальна відповідальність.} Міжнародний дискурс дизайн-освіти дедалі наполегливіше акцентує сталий дизайн та соціальну відповідальність як ключові компетентності~\cite{Maia2022105}. Специфікації стандарту B2 приділяють обмежену артикульовану увагу цим вимірам, хоча вони можуть адресуватися на рівні реалізації конкретних освітньо-професійних програм. У контексті ІКК ця лакуна означає недостатню увагу до етичних аспектів використання ШІ в дизайні, проблем упередженості алгоритмів та відповідального ставлення до генеративних технологій.

Окрім цих трьох лакун, виявлених у~\cite{Musiienko2025DesignEducation}, у контексті нашого дослідження необхідно наголосити на четвертій суттєвій лакуні --- \emph{цифрових/ШІ-навичках}. Стандарт B2 адресує ІКТ-навички (ЗК8) та використання програмного забезпечення для дизайну (ФК9) на загальному рівні, проте не специфікує компетентності, пов'язані зі штучним інтелектом: ШІ-грамотність, промпт-інженерію, критичне оцінювання ШІ-генерованого контенту, розуміння принципів роботи генеративних моделей. Як засвідчують дослідження~\cite{Song2025401, Hwang2025239, Chen2025}, інтеграція ШІ-технологій у дизайн-освіту є однією з найбільш актуальних тенденцій сучасності, а відсутність відповідних компетентностей у стандарті створює значний розрив між нормативними вимогами та реальними потребами ринку праці.

\textbf{Студійна педагогіка як основа формування дизайнерських компетентностей.} Синтез літератури з дизайн-освіти~\cite{Musiienko2025DesignEducation} підтверджує, що студійне навчання (studio-based learning) є педагогічним підходом з найсильнішим теоретичним обґрунтуванням для формування дизайнерських компетентностей. Студія залишається характерним педагогічним середовищем дизайн-освіти, що вирізняється інтеграцією індивідуальної творчої роботи, колаборативного навчання та експертного наставництва~\cite{Goldschmidt2010285}.

З позицій теорії рефлексивної практики Д.~Шона дизайн-студія створює умови для \emph{рефлексії-в-дії} (reflection-in-action) --- здатності, визначеної Шоном як центральної для розвитку професійної експертності~\cite{Schon1983}. Ітеративний цикл проєктування, отримання зворотного зв'язку та перегляду дозволяє студентам розвивати неявне знання (tacit knowledge) через досвідне залучення, а не лише через абстрактне навчання~\cite{Mewburn2012363}. Принцип <<навчання через діяльність>> є фундаментальним для студійної педагогіки: студенти розвивають компетентності не через засвоєння деконтекстуалізованих правил, а через занурення у реальний дизайнерський процес зі всією його складністю, невизначеністю та відкритістю.

З позицій дизайн-когніції студійне середовище підтримує коеволюцію проблеми та рішення, характерну для експертного дизайнерського мислення~\cite{Cross2011}. На відміну від інструкціоністських підходів, що розділяють визначення проблеми та генерування рішення, студійні проєкти уможливлюють одночасний розвиток обох процесів, формуючи когнітивні патерни, узгоджені з професійною практикою~\cite{Lawson2005}. \citet{Cross2018691} синтезує п'ятдесят років досліджень дизайну, артикулюючи концепт <<дизайнерських способів пізнання>> (designerly ways of knowing), що встановлює дизайн як легітимну дисципліну з характерними епістемологічними ознаками. Емпіричні дослідження засвідчують, що експертні дизайнери демонструють характерні когнітивні патерни, зокрема абдуктивне мислення, орієнтованість на рішення та ітеративне вдосконалення через прототипування~\cite{Kokotovich200849, Atman2019125}.

З позицій конструктивного узгодження (constructive alignment) Дж.~Біґґса оцінювання в студійному середовищі через критику проєктів (desk critique) уможливлює безпосередню оцінку дизайнерських компетентностей, продемонстрованих через дизайнерські артефакти та документацію процесу~\cite{Biggs1996}. Така узгодженість між навчальною діяльністю (проєктування) та оцінюванням (аналіз проєктів) підтримує когерентний розвиток компетентностей.

Для формування ІКК студійна педагогіка має особливе значення, оскільки студійне середовище природно інтегрує усі чотири компоненти ІКК: інформаційно-аналітичний (пошук референсів, аналіз контексту, дослідження користувачів), комунікативний (презентація проєктів, обговорення з наставником, командна робота), технологічний (використання цифрових інструментів проєктування), рефлексивний (аналіз власного процесу, реагування на критику, вдосконалення рішень). Саме інтегративний характер студійної педагогіки робить її найбільш адекватним середовищем для формування ІКК як цілісної характеристики.

\textbf{Проєктне навчання та доповнювальні педагогічні підходи.} Хоча студійне навчання забезпечує фундамент, аналіз~\cite{Musiienko2025DesignEducation} ідентифікує комплементарні підходи, що посилюють специфічні компетентнісні виміри.

\emph{Проєктне навчання з автентичними завданнями} розширює студійну педагогіку, розташовуючи дизайнерські виклики в реалістичних професійних контекстах~\cite{Kamalipour2022}. Автентичні проєкти розвивають компетентності клієнтської комунікації, управління обмеженнями та професійної відповідальності, які студійні вправи з штучними завданнями можуть недостатньо адресувати. Для формування комунікативного компонента ІКК проєктне навчання є особливо цінним, оскільки занурює студентів у реальні комунікативні ситуації з множинними стейкхолдерами. \citet{Tsai2023} засвідчує, що дизайнерське мислення в поєднанні з конструктивістським навчанням підвищує мотивацію та здатність до розв'язання складних проблем.

\emph{Структурована рефлексивна практика} адресує метакогнітивні виміри, ідентифіковані як недорозвинені в багатьох дизайнерських програмах. Артикульоване навчання рефлексії, підтримуване такими інструментами як дизайнерські журнали, портфоліо процесу та систематична самооцінка, розвиває компетентності критичної рефлексії, що відрізняють експертних практиків~\cite{Hameed2023310}. \citet{Hameed2023310} демонструє, що цифрова рефлексивна практика в дизайн-освіті текстилю посилює креативність через уможливлення повторного відвідування та реконструкції дизайнерської роботи. Це безпосередньо кореспондує з рефлексивним компонентом ІКК та є особливо актуальним в контексті рефлексії взаємодії з ШІ-інструментами.

\emph{Міждисциплінарні колаборативні проєкти} розвивають компетентності колаборації, дедалі більш затребувані у професійній практиці. Проєкти, що потребують залучення студентів з інших дисциплін --- інженерії, бізнесу, соціальних наук --- формують компетентності крос-дисциплінарної комунікації та інтеграції~\cite{Chon2019187}. \citet{Chon2019187} розрізняє дизайнерське мислення (процедурне знання) та дизайнерське пізнання (концептуальне розуміння), стверджуючи, що дизайн-освіта має культивувати обидва виміри через ретельно структурований педагогічний досвід.

Узагальнення зв'язку педагогічних підходів із формуванням конкретних компетентностей та компонентів ІКК представлено в табл.~\ref{tab:pedagogy-competency-ikk}.

\begin{table}[htbp]
\caption{Педагогічні підходи, компетентності та компоненти ІКК у дизайн-освіті}
\label{tab:pedagogy-competency-ikk}
\centering\small
\begin{tabularx}{\textwidth}{@{}p{2.8cm}Xp{2.5cm}p{2.5cm}@{}}
\toprule
\textbf{Педагогічний підхід} & \textbf{Основні компетентності} & \textbf{Компоненти ІКК} & \textbf{Теоретична основа} \\
\midrule
Студійне навчання & ХПД; навички дизайн-процесу; естетичне судження; технічне виконання & Усі чотири компоненти (інтегративно) & Рефлексивна практика \cite{Schon1983}; дизайн-когніція \cite{Cross2011} \\
\addlinespace
Проєктне навчання & Розв'язання проблем; професійна комунікація; управління обмеженнями & Комунікативний, технологічний & Досвідне навчання \cite{Kolb1984}; конструктивне узгодження \cite{Biggs1996} \\
\addlinespace
Структурована рефлексія & Критичне мислення; метакогніція; самооцінка; професійна ідентичність & Рефлексивний, інформаційно-аналітичний & Рефлексивна практика \cite{Schon1983, Hameed2023310} \\
\addlinespace
Колаборативні проєкти & Командна робота; міждисциплінарна комунікація; інтеграційні навички & Комунікативний & Соціальний конструктивізм; ситуативне навчання \\
\addlinespace
Дизайн-дослідження & Аналітичні здатності; доказовий дизайн; розуміння користувача & Інформаційно-аналітичний, рефлексивний & Дизайн-когніція \cite{Cross2011}; методологія дослідження \\
\bottomrule
\end{tabularx}
\end{table}

\textbf{Концепція художньо-проєктної діяльності як інтегративного принципу.} Концепт <<художньо-проєктної діяльності>> (ХПД), що є центральним для стандарту B2, заслуговує на окремий аналіз, оскільки він відображає специфічну українську традицію дизайн-освіти та водночас містить значний потенціал для інтеграції з сучасними європейськими підходами~\cite{Musiienko2025DesignEducation}.

Укорінений у радянській традиції, що інтегрувала підготовку з образотворчого мистецтва та прикладного дизайну, концепт ХПД організовує фахові компетентності у спосіб, що поєднує естетичні та технічні виміри. Замість розмежування між художньою творчістю та технічним проєктуванням, ХПД позиціонує їх як взаємопов'язані аспекти єдиного професійного процесу~\cite{Fursa2014}. Ця інтерпретація узгоджується з аргументом \citet{Cross2011} про те, що дизайн становить окремий когнітивний домен, який не зводиться ані до мистецтва, ані до інженерії.

Для формування ІКК концепт ХПД є продуктивним, оскільки він природно інтегрує інформаційну діяльність (аналіз контексту, пошук прототипів), комунікативну діяльність (презентація ідей, діалог з замовником), технологічну діяльність (використання інструментів проєктування) та рефлексивну діяльність (оцінювання естетичної якості, критичний аналіз рішень) у єдиний процесуальний цикл. Однак в умовах цифрової трансформації зміст кожного із цих видів діяльності суттєво змінюється: інформаційна діяльність включає взаємодію з ШІ-системами, комунікативна --- діалог з генеративними моделями через промпти, технологічна --- використання ШІ-інструментів генерації та редагування, рефлексивна --- критичне оцінювання ШІ-результатів та етичний аналіз їхнього використання.

\textbf{Виклики європейської інтеграції для дизайн-освіти.} Інтеграція української дизайн-освіти до Європейського простору вищої освіти, ініційована приєднанням до Болонського процесу у 2005~р. на Берґенській конференції~\cite{Bologna1999}, породжує низку суттєвих напружень, що потребують навігації~\cite{Musiienko2025DesignEducation, Castillo2011230}.

\emph{Стандартизація vs.\ креативність.} Європейські кваліфікаційні рамки наголошують на стандартизованих результатах навчання, придатних для міжнаціонального порівняння та визнання, однак дизайн-освіта за своєю природою цінує оригінальність та креативну дивергенцію, що чинять опір стандартизації~\cite{Weldon2009168}. Українська дизайн-освіта має відповідати вимогам рамок, не обмежуючи при цьому креативне дослідження, що є сутнісно необхідним для дизайнерського навчання. Як зазначає \citet{Brockmann200899}, перформансні результати навчання містять іманентні обмеження при застосуванні без урахування умов та контекстів, у яких демонструється компетентність.

\emph{Артикульоване vs.\ неявне знання.} Компетентнісні рамки вимагають артикульованого формулювання результатів навчання заради прозорості та оцінюваності, проте дизайнерська експертність значною мірою включає неявне знання --- тілесні здатності, естетичні чутливості та інтуїтивні судження, --- яке чинить опір повній артикуляції~\cite{Schon1983}. Розроблення підходів до оцінювання, які фіксують неявні компетентності в межах вимог рамок, залишається методологічним викликом. Це напруження особливо гостро проявляється у контексті ІКК, де значна частина інформаційно-комунікативної діяльності дизайнера базується на неявному знанні: інтуїтивному відборі візуальних рішень, <<відчутті>> доцільності комунікативного каналу, естетичній оцінці цифрового контенту.

\emph{Міжнародна гармонізація vs.\ національні традиції.} Гармонізація Болонського процесу сприяє мобільності та визнанню, але може також тиснути у напрямку конвергенції з домінантними європейськими моделями. Українська дизайн-освіта має характерні традиції --- інтеграцію підготовки з образотворчого мистецтва та прикладного дизайну, наголос на рисунку та ремісничих навичках, специфічну естетичну орієнтацію --- які заслуговують на збереження навіть у процесі європейської інтеграції~\cite{Fursa2014}.

\emph{Ресурсні обмеження.} Ефективна дизайн-освіта потребує суттєвих ресурсів: спеціалізованих студійних просторів, низького співвідношення студентів до викладачів для індивідуальної критики, матеріалів та обладнання для прототипування. Очікування європейської інтеграції можуть перевищувати ресурсну доступність в українських закладах, особливо з урахуванням триваючих дестабілізуючих факторів~\cite{Musiienko2025DesignEducation}. Для формування ІКК ця проблема набуває додаткового виміру: впровадження ШІ-інструментів у навчальний процес потребує відповідної технічної інфраструктури, ліцензій на програмне забезпечення та підвищення кваліфікації викладачів.

\textbf{Можливості європейської інтеграції.} Поряд із викликами, узгодження з європейськими кваліфікаційними структурами відкриває значні можливості для розвитку дизайн-освіти в Україні~\cite{Musiienko2025DesignEducation}.

Посилення мобільності та визнання є першорядною перевагою: узгодження сприяє руху студентів та викладачів у межах Європейського простору вищої освіти, забезпечуючи доступ до програм обміну, збагачуючи освітній досвід та формуючи професійні мережі, що перетинають національні кордони~\cite{Bohlinger2012279}. Для комунікативного компонента ІКК мобільність має особливе значення, оскільки занурює студентів у мультикультурні та мультимовні професійні контексти, що сприяє формуванню крос-культурної комунікативної компетентності.

Розвиток систем забезпечення якості є ще одним важливим напрямом. Європейські механізми забезпечення якості надають рамки для систематичної оцінки та вдосконалення програм~\cite{Srivastava2024305}. Залучення до цих механізмів може посилити дизайн-освіту через структуровану самооцінку та зовнішній перегляд програм.

Доступ до ресурсів розроблення навчальних програм є додатковою перевагою. Європейські мережі надають доступ до експертизи з розроблення програм, педагогічних інновацій та можливостей для колаборації. Українські заклади можуть використовувати європейський досвід, водночас привносячи унікальні перспективи з власних традицій, створюючи реципрокний обмін знаннями~\cite{Gonzalez2003}.

Покращення підготовки до ринку праці є фундаментальною метою інтеграції. Узгодження з європейськими компетентнісними рамками може посилити конкурентоспроможність українських випускників-дизайнерів на європейському ринку праці, розширюючи професійні можливості для окремих випускників та сприяючи інтеграції українських креативних індустрій до європейських ринків.

\textbf{Концептуальна рамка адаптивної інтеграції.} На основі проведеного аналізу в~\cite{Musiienko2025DesignEducation} було запропоновано концептуальну рамку <<адаптивної інтеграції>> як альтернативу поляризованим позиціям, що або некритично приймають гармонізацію, або захисно протистоять європейському залученню. Адаптивна інтеграція визнає легітимні інтереси, що обслуговуються європейськими рамками, --- мобільність, визнання, забезпечення якості, --- водночас наполягаючи на просторі для характерних національних та дисциплінарних традицій~\cite{Bohlinger2012279}.

Цей підхід спирається на теоретичні засади, що підтримують позицію щодо інкрементальної адаптації замість оптового запозичення: характерні ознаки дизайн-освіти --- центральність неявного знання, цінність креативної оригінальності, інтеграція естетичного та технічного судження --- створюють напруження з рамками, розробленими переважно для дисциплін із більш кодифікованим знанням. Адаптивна інтеграція навігує ці напруження через селективне прийняття, а не оптову відповідність~\cite{Musiienko2025DesignEducation}.

На рис.~\ref{fig:adaptive-integration} представлено концептуальну рамку адаптивної інтеграції, що позиціонує чотири стратегічні напрями навколо центрального принципу гармонізації із збереженням характерних сильних сторін.

\begin{figure}[htbp]
\centering
\begin{tikzpicture}[
  node distance=1.5cm,
  principle/.style={rectangle, draw=black, thick, rounded corners, minimum width=5.5cm, minimum height=1.8cm, align=center, fill=blue!10},
  strategy/.style={rectangle, draw=gray!60, rounded corners, minimum width=4.5cm, minimum height=1.2cm, align=center, fill=gray!5, font=\small},
  ikk/.style={rectangle, draw=black!60, rounded corners, minimum width=4.5cm, minimum height=0.8cm, align=center, fill=yellow!10, font=\footnotesize},
  arrow/.style={->, thick, >=stealth, draw=black!60}
]

% Core principle
\node[principle] (core) at (0,0) {\textbf{Адаптивна інтеграція}\\[4pt]Гармонізація із збереженням\\характерних сильних сторін};

% Four strategic directions
\node[strategy] (s1) at (-5,3) {Компетентнісне\\узгодження\\з гнучкістю};
\node[strategy] (s2) at (5,3) {Педагогічні\\інновації\\в межах традиції};
\node[strategy] (s3) at (-5,-3) {Забезпечення\\якості через\\діалог};
\node[strategy] (s4) at (5,-3) {Розвиток\\ресурсів для\\сталості};

% IKK connections
\node[ikk] (ikk1) at (-5,4.8) {ІКК: інформ.-аналіт.,};
\node[ikk] (ikk2) at (5,4.8) {ІКК: технолог., рефлекс.};
\node[ikk] (ikk3) at (-5,-4.8) {ІКК: рефлексивний};
\node[ikk] (ikk4) at (5,-4.8) {ІКК: технологічний};

% Arrows from core to strategies
\draw[arrow] (core) -- (s1);
\draw[arrow] (core) -- (s2);
\draw[arrow] (core) -- (s3);
\draw[arrow] (core) -- (s4);

% Arrows from strategies to IKK
\draw[arrow, gray!50, dashed] (s1) -- (ikk1);
\draw[arrow, gray!50, dashed] (s2) -- (ikk2);
\draw[arrow, gray!50, dashed] (s3) -- (ikk3);
\draw[arrow, gray!50, dashed] (s4) -- (ikk4);

\end{tikzpicture}
\caption{Концептуальна рамка адаптивної інтеграції дизайн-освіти в контексті європейської гармонізації (адаптовано за~\cite{Musiienko2025DesignEducation}): стратегічні напрями та зв'язок із компонентами ІКК}
\label{fig:adaptive-integration}
\end{figure}

Чотири стратегічні напрями операціоналізують принцип адаптивної інтеграції~\cite{Musiienko2025DesignEducation}:

\emph{Компетентнісне узгодження з гнучкістю} передбачає забезпечення чіткого зіставлення компетентностей B2 з ЄРК при збереженні простору для адаптації на рівні програм та характерних акцентів. У контексті ІКК це означає збереження можливості розширювати інформаційно-аналітичний та комунікативний компоненти відповідно до потреб конкретних спеціалізацій дизайну.

\emph{Педагогічні інновації в межах традиції} передбачають запозичення ефективних педагогічних підходів з міжнародної практики, спираючись на сильні сторони українських традицій дизайн-освіти. Для формування ІКК це означає інтеграцію ШІ-інструментів у традиційне студійне навчання, а не заміну студійної педагогіки технологічно-орієнтованими підходами. Технологічний та рефлексивний компоненти ІКК розвиваються через доповнення, а не заміщення традиційних практик.

\emph{Забезпечення якості через діалог} передбачає залучення до європейських систем забезпечення якості як можливості для розвитку, а не тягаря відповідності, з привнесенням українських перспектив до європейського дискурсу. Рефлексивний компонент ІКК безпосередньо сприяє цьому напряму, оскільки сформована здатність до критичної самооцінки є необхідною умовою продуктивної участі у процесах забезпечення якості.

\emph{Розвиток ресурсів для сталості} передбачає стратегічне адресування ресурсних потреб, використання європейських партнерств та можливостей фінансування при побудові сталого інституційного потенціалу. Для технологічного компонента ІКК це означає систематичне нарощування інфраструктури, необхідної для роботи з сучасними цифровими та ШІ-інструментами.

\textbf{Інтегративна теоретична рамка дизайн-освіти.} Проведений аналіз дозволяє сформулювати інтегративну теоретичну рамку, що позиціонує розвиток компетентностей дизайн-освіти на перетині трьох взаємодіючих систем~\cite{Musiienko2025DesignEducation}: системи специфікації компетентностей (формальні рамки та стандарти), педагогічної системи (підходи до навчання) та системи інтеграції (процеси європейської гармонізації). У контексті нашого дослідження ця рамка доповнюється четвертим виміром --- ІКК як інтегративною характеристикою, що пронизує всі три системи.

Система специфікації компетентностей встановлює формальні вимоги через такі рамки як ЄРК та національні стандарти, включаючи стандарт B2 <<Дизайн>>. Педагогічна система охоплює підходи до навчання --- студійне навчання, проєктну педагогіку, рефлексивну практику --- через які компетентності формуються. Система інтеграції репрезентує ширший політичний контекст європейської гармонізації, механізмів забезпечення якості та процесів визнання.

Критично важливо, що рамка акцентує двосторонні зв'язки між системами: специфікації компетентностей впливають на педагогічні вибори, але педагогічні традиції також формують спосіб інтерпретації та операціоналізації компетентностей; вимоги інтеграції створюють тиск до стандартизації, проте характерна природа дизайн-освіти генерує напруження зі стандартизованими підходами. Ці динамічні взаємодії конституюють аналітичний фокус нашого дослідження.

На рис.~\ref{fig:three-systems-ikk} представлено інтегративну теоретичну рамку розвитку компетентностей дизайн-освіти, доповнену виміром ІКК.

\begin{figure}[htbp]
\centering
\resizebox{\linewidth}{!}{
\begin{tikzpicture}[
  corenode/.style={circle, draw=black!80, very thick, minimum size=3.2cm, align=center, fill=yellow!15},
  systembox/.style={rectangle, draw=black, thick, rounded corners=8pt, minimum width=4.5cm, minimum height=2.2cm, align=center, fill=blue!8},
  theorybox/.style={rectangle, draw=gray!60, rounded corners=4pt, minimum width=3.8cm, minimum height=0.7cm, align=center, fill=white, font=\footnotesize},
  ikkbox/.style={rectangle, draw=black!60, very thick, rounded corners=6pt, minimum width=4cm, minimum height=1.2cm, align=center, fill=yellow!8, font=\small},
  processarrow/.style={->, very thick, >=stealth, draw=blue!60},
  feedbackarrow/.style={<->, thick, >=stealth, draw=gray!50, dashed},
  ikkconnect/.style={<->, thick, >=stealth, draw=black!40, densely dotted}
]

% Central core
\node[corenode] (core) at (0,0) {\textbf{Формування}\\[3pt]\textbf{компетентностей}\\[3pt]\textbf{дизайнера}};

% Three system nodes
\node[systembox] (competency) at (0,5) {\textbf{Система}\\[2pt]\textbf{специфікації}\\[2pt]\textbf{компетентностей}};
\node[systembox] (pedagogy) at (-5,-3) {\textbf{Педагогічна}\\[2pt]\textbf{система}};
\node[systembox] (integration) at (5,-3) {\textbf{Система}\\[2pt]\textbf{інтеграції}};

% Theory boxes for Competency System
\node[theorybox] (eqf) at (-2,6.8) {ЄРК (EQF)};
\node[theorybox] (b2) at (2,6.8) {Стандарт B2 022};

% Theory boxes for Pedagogical System
\node[theorybox] (schon) at (-9.5,-1.4) {Рефлексивна практика\\(Шон)};
\node[theorybox] (biggs) at (-9.5,-2.8) {Конструктивне узгодження\\(Біґґс)};
\node[theorybox] (design) at (-9.5,-4.2) {Дизайн-когніція\\(Кросс, Лоусон)};

% Theory boxes for Integration System
\node[theorybox] (bologna) at (9.5,-1.4) {Болонський процес};
\node[theorybox] (qa) at (9.5,-2.8) {Забезпечення якості};
\node[theorybox] (mobility) at (9.5,-4.2) {Мобільність та визнання};

% IKK dimension at bottom
\node[ikkbox] (ikk) at (0,-6) {\textbf{ІКК дизайнера}\\інформ. $\cdot$ комунік. $\cdot$ технол. $\cdot$ рефлекс.};

% Main arrows from systems to core
\draw[processarrow] (competency) -- node[left, font=\scriptsize\itshape, text=gray!60] {результати} (core);
\draw[processarrow] (pedagogy) -- node[below left, font=\scriptsize\itshape, text=gray!60] {методи} (core);
\draw[processarrow] (integration) -- node[below right, font=\scriptsize\itshape, text=gray!60] {контекст} (core);

% Theory box connections
\draw[gray!40, thick] (eqf) -- (competency);
\draw[gray!40, thick] (b2) -- (competency);
\draw[gray!40, thick] (schon) -- (pedagogy);
\draw[gray!40, thick] (biggs) -- (pedagogy);
\draw[gray!40, thick] (design) -- (pedagogy);
\draw[gray!40, thick] (bologna) -- (integration);
\draw[gray!40, thick] (qa) -- (integration);
\draw[gray!40, thick] (mobility) -- (integration);

% Inter-system feedback loops
\draw[feedbackarrow] (competency.west) to[bend right=20] node[above left, font=\scriptsize\itshape, text=gray!60] {узгодження} (pedagogy.north);
\draw[feedbackarrow] (competency.east) to[bend left=20] node[above right, font=\scriptsize\itshape, text=gray!60] {гармонізація} (integration.north);
\draw[feedbackarrow] (pedagogy.east) to[bend right=15] node[below, font=\scriptsize\itshape, text=gray!60] {реалізація} (integration.west);

% IKK connections
\draw[ikkconnect] (core) -- (ikk);
\draw[ikkconnect] (pedagogy.south) to[bend right=20] (ikk.west);
\draw[ikkconnect] (integration.south) to[bend left=20] (ikk.east);

\end{tikzpicture}
}
\caption{Інтегративна теоретична рамка розвитку компетентностей дизайн-освіти в контексті європейської інтеграції та формування ІКК (адаптовано за~\cite{Musiienko2025DesignEducation}). Суцільні стрілки позначають прямий вплив; штрихові --- двосторонній зворотний зв'язок; пунктирні --- зв'язок із ІКК}
\label{fig:three-systems-ikk}
\end{figure}

Як видно з рис.~\ref{fig:three-systems-ikk}, ІКК позиціонується як наскрізна характеристика, що формується під впливом усіх трьох систем: специфікація компетентностей визначає нормативні вимоги до ІКК (через ЗК7, ЗК8, ФК9 стандарту B2); педагогічна система забезпечує методичний інструментарій формування ІКК (через студійне, проєктне навчання та рефлексивну практику); система інтеграції створює контекст, у якому ІКК набуває додаткових вимірів (крос-культурна комунікація, міжнародні стандарти цифрової грамотності, європейські вимоги до якості).

\textbf{Особливості дизайн-освіти в пострадянському контексті.} Дизайн-освіта в пострадянських країнах розвивалася в специфічних історичних обставинах, виростаючи з традицій, що акцентували технічну підготовку в спеціалізованих академіях, відокремлених від закладів образотворчого мистецтва~\cite{Fursa2014}. Радянська модель надавала пріоритет промисловому застосуванню та колективним виробничим цінностям, при цьому естетичні міркування підпорядковувалися функціональним та ідеологічним вимогам. Ця спадщина продовжує впливати на інституційні структури, експертність викладацького складу та педагогічні традиції в сучасній дизайн-освіті регіону.

Приєднання України до Болонського процесу у 2005~р. ініціювало систематичну реформу вищої освіти, включаючи дизайнерські програми. Затвердження стандарту B2 <<Дизайн>> у 2018~р. є значною віхою, що артикулює компетентнісні вимоги, узгоджені з європейськими рамками, адресуючи водночас потреби національного ринку праці~\cite{MON2018}. Стандарт визначає як загальні компетентності (критичне мислення, комунікація, командна робота), так і фахові компетентності, специфічні для дизайнерської практики (художньо-проєктна діяльність, естетичний аналіз, технічна реалізація).

Проте кілька викликів ускладнюють процес реформування. Розвиток викладацького складу є суттєвою проблемою: значна частина дизайн-педагогів була підготовлена в межах радянських традицій і може не мати достатньої обізнаності з компетентнісними підходами та сучасними міжнародними дизайнерськими практиками. Ресурсні обмеження лімітують доступ до сучасних технологій та можливостей міжнародного обміну. Триваючий конфлікт додатково порушив освітню безперервність, призвівши до переміщення студентів та викладачів.

Незважаючи на ці виклики, українська дизайн-освіта демонструє значний динамізм. Інтеграція з європейськими мережами розширила можливості для мобільності студентів та викладачів. Реформування навчальних програм увело сучасні підходи, включаючи методології дизайн-мислення та цифрові технології дизайну. Стандарт B2 забезпечує рамку для систематичного розвитку компетентностей, хоча реалізація варіюється між закладами~\cite{Musiienko2025DesignEducation}.

Для формування ІКК контекст пострадянської трансформації має подвійне значення. З одного боку, сильна традиція фундаментальної художньої підготовки забезпечує міцну основу для рефлексивного компонента ІКК (розвинена естетична чутливість, здатність до критичного аналізу візуальної інформації). З іншого --- відносна ізоляція від глобальних тенденцій цифровізації дизайну та запізніле впровадження ШІ-технологій створюють дефіцити в технологічному компоненті ІКК, що потребують цілеспрямованого подолання.

\textbf{Технологічна трансформація та ІКК.} Інтеграція нових технологій у дизайн-освіту є одним з найбільш суттєвих сучасних викликів для європейських закладів~\cite{Chen2025}. \citet{Song2025401} документує стрімкий перехід від загального ШІ до генеративного ШІ у дизайн-освіті, де генерація зображень стає найбільш поширеною технологією. Ця технологічна трансформація потребує адаптації навчальних програм, водночас порушуючи питання про природу дизайнерської креативності та про компетентності, що відрізняють людського дизайнера від машини.

\citet{Hwang2025239} стверджує, що дизайн-освіта в еру генеративного ШІ потребує фундаментально інших освітніх цілей, наголошуючи на ШІ-візуальній грамотності та ролі дизайнерів як створювачів контенту, що поєднують науку, мистецтво та технологію. Це узгоджується з виявленим у підрозділі~\ref{subsec:1.2} зсувом парадигми <<від виробника до куратора>>, що суттєво трансформує зміст усіх компонентів ІКК.

\citet{Melker20251871} пропонує концептуальний підхід до посилення дивергентного та конвергентного мислення студентів-дизайнерів у процесах ідеації через використання ШІ. \citet{Chandrasekera2025358} досліджує інструменти генеративного ШІ для різних стилів навчання у дизайн-освіті. \citet{Yusuf2026} порівнює індивідуальну та парну підтримку в еру ШІ, досліджуючи когнітивні диспозиції студентів у дизайн-освіті. Ці дослідження засвідчують зростаючу інтеграцію ШІ-технологій у педагогічну практику дизайн-освіти та необхідність відповідного оновлення компетентнісних вимог.

Для технологічного компонента ІКК зазначені тенденції означають необхідність включення нових елементів: ШІ-грамотності (розуміння принципів роботи генеративних моделей), промпт-інженерії (здатність формулювати ефективні запити до ШІ-систем), критичного оцінювання ШІ-генерованого контенту (розрізнення якісних та неякісних результатів, виявлення упередженостей), етичного використання генеративних технологій (проблеми авторського права, плагіату, прозорості). Для рефлексивного компонента ІКК технологічна трансформація додає новий вимір --- здатність критично осмислювати власну взаємодію з ШІ-інструментами, усвідомлювати межі автоматизованих рішень та підтримувати авторську позицію.

\textbf{Досвідне навчання та формування ІКК.} Окремо зазначимо роль досвідного навчання (experiential learning) у формуванні ІКК дизайнерів. Модель досвідного навчання Д.~Колба~\cite{Kolb1984}, що включає чотири фази --- конкретний досвід, рефлексивне спостереження, абстрактну концептуалізацію та активне експериментування --- добре узгоджується з циклічним характером дизайнерського процесу та природно інтегрує компоненти ІКК. \citet{Duchatelet2024400} у систематичному огляді середовищ досвідного навчання у вищій освіті демонструє зв'язок між характеристиками середовищ досвідного навчання та формуванням генеричних результатів навчання, що включають критичне мислення, комунікативні навички та здатність до самоуправління --- тобто ключові складники ІКК.

У контексті формування ІКК дизайнерів модель Колба може бути інтерпретована наступним чином: конкретний досвід реалізується через безпосередню роботу з цифровими та ШІ-інструментами у студійному середовищі (технологічний компонент); рефлексивне спостереження --- через аналіз результатів взаємодії з технологіями та оцінювання ефективності комунікативних стратегій (рефлексивний компонент); абстрактна концептуалізація --- через формування узагальнених моделей інформаційної поведінки та принципів ефективної цифрової комунікації (інформаційно-аналітичний компонент); активне експериментування --- через апробацію нових підходів до використання ШІ-інструментів та комунікативних стратегій у реальних дизайнерських проєктах (комунікативний та технологічний компоненти).

\textbf{Проблема оцінювання дизайнерських компетентностей.} Окремим аспектом, що потребує уваги в контексті європейської інтеграції, є проблема оцінювання дизайнерських компетентностей у компетентнісних рамках. Напруження між стандартизацією та креативністю створює особливі виклики для дизайн-освіти~\cite{Weldon2009168}. \citet{AlvarezIcaza2024} адресує оцінювання креативної поведінки студентів вищої освіти, розрізняючи об'єктивні профілі діяльності та суб'єктивне самосприйняття. Їхнє дослідження зі студентами програм дизайну, архітектури та комунікацій виявляє, що студенти-дизайнери демонструють характерні патерни креативної поведінки, хоча значна індивідуальна варіація існує.

\citet{Ardington2017157} аналізує дискурс оцінювання в дизайн-студії, виявляючи напруження між потребою в артикульованому керівництві та збереженням креативної свободи, сутнісно необхідної для дизайнерського навчання. Суб'єктивний характер дизайнерського оцінювання, опосередкованого дискурсом, заснованим на неявних дисциплінарних розуміннях, створює постійні виклики для компетентнісних рамок, що потребують артикульованого формулювання результатів навчання.

Для оцінювання ІКК дизайнерів ці виклики набувають додаткового виміру. Якщо оцінювання традиційних дизайнерських компетентностей (естетичне судження, технічна майстерність) вже є складним завданням, то оцінювання ІКК потребує врахування ще більш складних параметрів: якості інформаційної поведінки, ефективності цифрової комунікації, релевантності вибору технологічних інструментів, глибини рефлексії щодо взаємодії з ШІ-системами. Це зумовлює необхідність розроблення спеціалізованих оцінювальних інструментів, що є одним із завдань подальшого дослідження.

\textbf{Аналітичний синтез підрозділу.} Проведений аналіз дозволяє сформулювати кілька ключових висновків, що мають значення для подальшого дослідження формування ІКК майбутніх дизайнерів.

По-перше, стандарт B2 <<Дизайн>> демонструє сильну структурну відповідність архітектурі ЄРК, що забезпечує формальні вимоги європейської інтеграції --- визнання кваліфікацій, трансферт кредитів, порівнюваність кваліфікацій. Водночас існують суттєві концептуальні варіації, зокрема щодо унікального українського концепту <<художньо-проєктної діяльності>>, який інтегрує естетичні та технічні виміри у спосіб, що не має прямих еквівалентів у західноєвропейських рамках. Цей концепт слід розглядати не як дефіцит, що потребує корекції, а як характерну сильну сторону --- інтегративне розуміння дизайнерської компетентності.

По-друге, систематичне зіставлення компетентностей стандарту B2 з компонентами ІКК засвідчує, що стандарт імпліцитно визначає структурні елементи ІКК, проте не артикулює їх як цілісну інтегративну характеристику. Ідентифіковано чотири компетентнісні лакуни: критична рефлексія, міждисциплінарна колаборація, сталий розвиток та цифрові/ШІ-навички, --- що потребують адресування для забезпечення повноцінного формування ІКК.

По-третє, студійне навчання, підтримуване проєктною педагогікою та структурованою рефлексивною практикою, становить найбільш теоретично обґрунтований підхід до формування дизайнерських компетентностей та, відповідно, ІКК. Інтегративний характер студійної педагогіки робить її найбільш адекватним середовищем для формування ІКК як цілісної характеристики, оскільки студійне середовище природно залучає усі чотири компоненти ІКК.

По-четверте, концепція адаптивної інтеграції~\cite{Musiienko2025DesignEducation} пропонує продуктивну альтернативу оптовому запозиченню європейських моделей, наголошуючи на інкрементальній адаптації, що зберігає сильні сторони національної традиції, водночас адресуючи виявлені лакуни. Цей підхід є особливо релевантним для формування ІКК, де необхідно інтегрувати нові цифрові та ШІ-компетентності, не руйнуючи фундаментальних засад художньо-проєктної підготовки.

Ці висновки утворюють нормативно-аналітичну основу для подальшого розроблення теоретичної рамки формування ІКК, що буде здійснено в наступному підрозділі~\ref{subsec:1.4} через інтеграцію моделі TPACK, студійної педагогіки та компетентнісного підходу.
